\documentclass{article}
\usepackage{graphicx} % Required for inserting images
\usepackage{amsmath}
\title{Lesson 4 Homework - Proof Writing}
\author{William Wang}
\date{September 2025}

\begin{document}

\maketitle

\section{}
We have \(\frac{A}{B+Ce^{-kt}} = D\), we can first get rid of the fraction. 
\[\frac{A}{B+Ce^{-kt}} = D \implies A = DB+DCe^{-kt}\]
\[\implies \frac{A-DB}{DC} = e^{-kt} \implies \ln(\frac{A-DB}{DC}) = -kt\]
\[t = -\frac{1}{k}\ln(\frac{A-DB}{DC})\]
\section{}
\[a+b\ln c = d \implies b\ln c = d-a \implies \ln c = \frac{d-a}{b}\]
\[\implies c = e^{\frac{d-a}{b}}\]
\section{}
Since we know the half time, remaining mass and time decayed, we have
\[4.2 = N_0(\frac{1}{2})^{12/4.3}\]
\[N_0 = 4.2/(\frac{1}{2})^{12/4.3} \approx 29.06\]
\section{}
Let $x$ be the initial amount and $t_{1/2}$ be the half life, then we have 
\[12.2 = x(\frac{1}{2})^{63/t_{1/2}} \text{ and } 4.1 = x(\frac{1}{2})^{107/t_{1/2}}\]

By dividing these two equations, we get
\[\frac{4.1}{12.2} = (\frac{1}{2})^{107/t_{1/2} - 63/t_{1/2}}\]
\[= (\frac{1}{2})^{44/t_{1/2}}\]
By taking the natural logarithm base of each side, we have
\[\ln \frac{4.1}{12.2} = \ln(\frac{1}{2})^{44/t_{1/2}} = \frac{44}{t_{1/2}}\ln \frac{1}{2}\]
\[\implies t_{1/2} = \frac{44\ln(0.5)}{\ln(4.1/12.2)} \approx 27.97\]
Hence, we have \(x = \frac{12.2}{(1/2)^{63/27.97}} \approx 58.13 \). 


\end{document}
